\section{Agent Runtime y Agent Loop}

\subsection{Agent Runtime: el entorno de ejecución del Agent}

Agent Runtime es el «escenario» en el que se ejecuta el Agent: provee todo el soporte de entorno que necesita para funcionar:

\begin{itemize}
    \item \textbf{Workspace}: como el espacio de trabajo de OpenAI Codex
    \item \textbf{Inyección de archivos de arranque}: el prompt a nivel de sistema (System Prompt)
    \item \textbf{Tools y Skills integrados}: el conjunto de herramientas predefinidas y la documentación de skills
    \item \textbf{Archivos de configuración}: archivos de configuración en Markdown como agents.md, src.md, tools.md, entre otros
    \item \textbf{Sistema de Memory}: el mecanismo de memoria y persistencia
\end{itemize}

\subsection{Agent Loop: el ciclo de ejecución impulsado por mensajes}

Agent Loop describe cómo, una vez que un mensaje entra realmente al sistema, se procesa hasta generar finalmente una respuesta o una acción:

\begin{importantbox}{El flujo estándar del Agent Loop}
\begin{enumerate}
    \item \textbf{Intake}: se recibe el mensaje del usuario (puede requerir ponerse en cola)
    \item \textbf{Prompt Injection}: se inyecta una gran cantidad de prompts (instrucciones del sistema, Skills, información del entorno, etc.)
    \item \textbf{Model Inference}: se llama a un modelo grande remoto o local para realizar la inferencia
    \item \textbf{Function Call} (opcional): si el modelo decide invocar una herramienta, se ejecuta la Function correspondiente
    \item \textbf{Streaming Response}: la respuesta se emite en flujo (streaming)
    \item \textbf{Persistence}: se persiste localmente (se guarda el historial de la conversación)
\end{enumerate}
Este es el flujo estándar basado en el Agent Loop descrito en la documentación oficial de OpenAI Codex.
\end{importantbox}

\begin{knowledgebox}{Agent Runtime vs. Agent Loop}
\begin{itemize}
    \item \textbf{Agent Runtime}: el entorno de ejecución estático, que se prepara cuando el Agent arranca y provee todos los recursos necesarios para la ejecución
    \item \textbf{Agent Loop}: el ciclo de ejecución dinámico que, al entrar un mensaje, impulsa al Agent a ejecutar acciones y generar respuestas
\end{itemize}
Runtime es el escenario; Loop es el proceso de la actuación sobre ese escenario.
\end{knowledgebox}

\subsection{Resumen de la sección}

Agent Runtime provee el entorno de ejecución (workspace, archivos de configuración, herramientas integradas); Agent Loop es el ciclo de ejecución impulsado por mensajes (recepción → inyección → inferencia → invocación → respuesta → persistencia). Juntos conforman la arquitectura completa de un Modern Agent.


